%% Auto-generated LaTeX version of user_manual.md
\documentclass[11pt,a4paper]{article}
\usepackage[utf8]{inputenc}
\usepackage{fontawesome5}
\usepackage[T1]{fontenc}
\usepackage{lmodern}
\usepackage{graphicx}
\usepackage{geometry}
\usepackage{float}
\geometry{margin=1in}

% Disable paragraph indentation throughout the document
\setlength{\parindent}{0pt}

\begin{document}
\pagestyle{empty}

% Top Header: Company Info (Left) and Logo (Right)
\noindent
\begin{minipage}[c]{0.5\textwidth}
    \textbf{Mlnarik Organ, s.r.o.} \\
    Cerhenska 9 \\
    29001 Sokolec \\
    Czech Republic
\end{minipage}
\begin{minipage}[c]{0.5\textwidth}
    \raggedleft
    
\includegraphics[height=1.5cm]{logo.png} % Replace with your image file
\end{minipage}

% Place title directly below the logo (avoid starting a new page)
\vspace{3cm}

\begin{center}
    {\Huge \textbf{HALL--SCANNER64: User Manual}}\\
    \vspace{0.3cm}
    February 2026

    \noindent Applicable to: Firmware Version 1.0

\end{center}

\vspace{1cm}

\begin{center}
    \includegraphics[width=1\textwidth]{Hall-scanner64.png}
\end{center}

\vfill

\newpage

\section{Overview}

HALL--SCANNER64 reads up to 64 analog Hall sensor channels (8~\texttimes~MCP3008 ADC chips, 8 channels each) over SPI and generates MIDI NOTE ON/OFF messages based on calibrated trigger thresholds. The module is specifically tailored for keyboards made by Mlnarik Organ, featuring Hall sensors.

It supports both traditional 5-pin DIN MIDI and USB-MIDI output with automatic merging of incoming MIDI data. In addition to standard MIDI capabilities, the module introduces several advanced features such as a High-Speed Analog MIDI Bus and an easy-to-use Wi-Fi configuration interface accessible via a web browser.

\section{Features}
\begin{itemize}
    \item Powered via 5~V DC / 0.5~A
    \item Reads up to 64 Hall sensors
    \item Sends velocity-sensitive MIDI NOTE ON messages (touch response)
    \item MIDI IN (5-pin DIN)
    \item MIDI OUT (5-pin DIN) with automatic merging of MIDI IN and internal messages
    \item USB MIDI output (class-compliant)
    \item Optional High-Speed Analog MIDI Bus -- supports up to 4\texttimes{} the standard MIDI baud rate for expanded bandwidth
    \item Wi-Fi Configuration Interface -- configure module settings via a web browser by connecting to the module's built-in Wi-Fi access point
\end{itemize}

\section{Introduction to MIDI Device Connectivity and Latency Considerations}

In practical applications, it is common to use multiple MIDI modules in a setup, often combining devices from various manufacturers. There are two basic approaches to connecting input MIDI devices to an end device such as a computer or Digital Audio Workstation (DAW).


\begin{figure}[H]
    \centering
    \includegraphics[width=0.2\textwidth]{Analog_chaining.drawio.png}
    \caption{Chaining via the analog MIDI bus -- this typically involves using traditional 5-pin DIN connectors (though other connectors may also be used).}
\end{figure}

\begin{figure}[H]
    \centering
    \includegraphics[width=0.8\textwidth]{USB_star.drawio.png}
    \caption{Connecting in a "star" topology via USB.}
\end{figure}

\begin{center}
    \fbox{
    \begin{minipage}{0.95\linewidth}
        \faInfoCircle~\textbf{Note: Latency considerations} If you prefer chaining via the analog MIDI bus and use many MIDI devices, latency can become a significant issue. The transmission of a single standard MIDI message takes approximately 1\,ms. As MIDI data is passed through multiple devices in the chain, this latency can accumulate and become noticeable.
    \end{minipage}
    }
\end{center}

One of the key features of the HALL--SCANNER64 module is its optional High-Speed Analog MIDI Bus mode. When activated, MIDI messages are transmitted at four times the standard MIDI speed. However, please note that this high-speed mode is not compatible with third-party MIDI devices, as they are designed to operate at the standard MIDI baud rate.

\section{Setup and installation}

\begin{figure}[H]
    \centering
    \includegraphics[width=1\textwidth]{labeled-board.drawio.png}
    \caption{Overall layout of the board}
\end{figure}

\subsection{Power}
\begin{itemize}
  \item Voltage: 5~V DC (regulated)
  \item Current Consumption: 0.5~A (typical, when all 64 Hall sensors are active)
\end{itemize}

\begin{center}
    \fbox{
    \begin{minipage}{0.95\linewidth}
        \faExclamationTriangle~\textbf{Warning:}
        \begin{itemize}
            \item Use a regulated 5~V supply. Do not exceed 5.3~V. Reverse polarity can damage the device.
            \item Use a power supply with sufficient current capacity, especially if powering multiple modules. Always use a dedicated power input; do not rely on USB power alone. High current draw from the Hall sensors can cause voltage drop, which may shift trigger points.
        \end{itemize}
    \end{minipage}
    }
\end{center}

\subsection{Hall Sensor Inputs (keyboard connection)}

The board expects a dedicated Hall sensor expansion board to be connected — the scanner does not read passive mechanical switches directly. Do not attempt to connect a switch matrix; instead, use the Hall sensor board provided by Mlnarik Organ.

\begin{center}
    \fbox{
    \begin{minipage}{0.95\linewidth}
        \faExclamationTriangle~\textbf{Warning:}
        Because a high-speed SPI bus is used to read the sensors, the flat ribbon cable between this module and the Hall sensor board should not exceed 1 m in length. Longer cables can cause signal integrity issues and unreliable sensor readings.
    \end{minipage}
    }
\end{center}


\section{Analog MIDI Connection}

The image below shows how to connect to the analog MIDI bus using DIN-5 connectors. However, it is not necessary to use connectors between the modules — you can also use wires connected directly to the screw terminals. Shielded cables are recommended.

\begin{figure}[H]
    \centering
    \includegraphics[width=1\textwidth]{MIDI-schema.png}
    \caption{Connection analog MIDI bus}
\end{figure}


\section{Setting Up the Module}

You can conveniently set up the module using any web browser on your computer or mobile phone. The module can start in a special mode where it creates a Wi-Fi access point that you can connect to.

    \subsection{Enabling Wi-Fi Access Point Mode}
\begin{enumerate}
    \item Unplug the module from any power source (USB and external power).
    \item Press and hold the WiFi button.
    \item Continue holding the button while powering on the module (e.g., by connecting an external power source).
    \item The Status LED will light up steadily, indicating that the module is now in Wi-Fi access
    point mode.
    \item On your computer (with a Wi-Fi adapter) or smartphone, go to the Wi-Fi settings. You
    should see a new Wi-Fi network named "Hall-scanner64". Connect to it.

    \faInfoCircle~\textbf{Note}: While connected to this network, your device will lose access to the internet.
    \item Open a web browser (such as Google Chrome) and enter the following IP address in the address bar: 192.168.4.1 -- the fixed IP address of the MIDI module's web server. You should now see the settings page, as shown in the image below.
    \item Use the dropdown menus to configure all the necessary settings, then press the
    "Save Settings" button. If needed, you can restore the factory default values by
    pressing the "Reset to Defaults" button.
    \item Once you're done, disconnect the module from the power source. Your settings will be saved and automatically applied the next time the module starts.

    \faInfoCircle~\textbf{Note}: While in Wi-Fi access point mode, the module does not transmit MIDI data. MIDI communication is disabled during setup.
\end{enumerate}

\begin{figure}[H]
    \centering
    \includegraphics[width=0.8\textwidth]{app-settings-window.png}
    \caption{Settings page in Wi-Fi Access Point mode}
\end{figure}

The firmware provides a minimal HTTP-based settings page when the unit runs in Wi-Fi Access Point mode. The page allows:
\begin{itemize}
  \item Changing MIDI channel (range: 1--16)
    \item Changing base MIDI note (range: 0--95). The first sensor maps to this note; subsequent sensors increment the note number.
    \item Fast MIDI: toggle between the standard MIDI baud rate (31250) and a higher UART rate (115200).
  \item Running the Keyboard calibration routine (start/stop)
\end{itemize}

\subsection{Keyboard Calibration}
This function is intended for Mlnarik Organ keyboards that do not use regulation screws for setting individual key trigger points. The calibration routine automatically determines appropriate per-key trigger thresholds. During calibration, the user presses and releases each key to be recalibrated. The module samples the "released" and "pressed" voltage states and then computes the individual trigger thresholds and required hysteresis values.

When the calibration is started, a new dialog will appear on the web page. It says "Gently press and release each key to be calibrated for at least 1 second."

\begin{figure}[H]
    \centering
    \includegraphics[width=0.8\textwidth]{keyboard-calibration.png}
    \caption{Keyboard calibration dialog}
\end{figure}

The calibration routine recognizes which keys moved and recalibrates only those keys. You do not need to calibrate each key individually -- you can press and release multiple keys at the same time. You may also press and release any key multiple times.

\begin{center}
    \fbox{
    \begin{minipage}{0.95\linewidth}
        \faExclamationTriangle~\textbf{Warning:}
        \begin{itemize}
            \item Press and release the desired keys gently and slowly to avoid key vibration, especially when released.
            \item Press the keys for approximately 1 second. In calibration mode, the voltage signal from each sensor is highly filtered and averaged to capture stable "pressed" and "released" voltage levels.
            \item Use the same stable supply voltage during calibration as in normal operation. Voltage shifts between calibration and operation can result in incorrect trigger points.
        \end{itemize}
    \end{minipage}
    }
\end{center}

\section{Status LED behavior}
\begin{itemize}
  \item Startup: brief blink sequence
  \item MIDI activity: short blink when MIDI bytes are transmitted
  \item AP mode: LED remains steadily lit while the AP is active
\end{itemize}

\end{document}
